\section{Uniform harmonic coordinates}
\label{co:section}

We first obtain an open coordinate domain.  The estimates in this
section allow the derivatives of these coordinates to degenerate as
the scale decreases.  Their degeneration will be measured by positive
matrices in the next section.

\subsection{Local calculus and rescaling}

For locally Sobolev functions write
\(\Gamma(f,g)=\langle\nabla f,\nabla g\rangle\) and
\(\Gamma(f)=\Gamma(f,f)\).  Laplacians of gradient squares will be
understood as measures.  Products with such measures always use the
continuous representative of a Lipschitz function, or the capacity
representative of a Sobolev function when that is the appropriate
meaning of the pairing.

We write \(g\) for the canonical metric tensor.

\begin{lemma}[Local analytic facts]\label{co:local-calculus}
Let \((Y,d,\m)\) be a noncollapsed \(\RCD(-\kappa,n)\) space and let
\(\Omega\subset Y\) be open.  A locally bounded solution of
\(\Delta w=a\) in \(\Omega\), with \(a\) constant, is locally Lipschitz
and is a local test function.  In particular it has a local
\(L^2\) Hessian, \(\Gamma(w_1,w_2)\in W^{1,2}_{\rm loc}(\Omega)\)
for two such functions, and the Hessian product rule holds locally.
The measure
\begin{equation}\label{co:M-definition}
 M(w)=\frac12\Delta\Gamma(w)+\kappa\Gamma(w)\m
\end{equation}
is locally finite and nonnegative.  For every constant \(b\),
\begin{equation}\label{co:shifted-bochner}
 M(w)-(2b\Delta w-nb^2)\m
 \ge |\Hess w-bg|^2\m.
\end{equation}
On the real vector space of local constant-Laplacian functions,
the left side is a nonnegative quadratic form in the pair
\((w,b)\), with both entries scaled together.
All these statements hold on interior balls for weak solutions with
locally bounded \(L^2\) norm, after using the interior estimates to
obtain local boundedness.

On fixed interior scales, bounded Lipschitz cutoffs with bounded
Laplacian are available.  Under uniform local volume comparison and
a uniform rescaled curvature bound, their bounds, local Poincare
and Sobolev constants, and interior harmonic and Poisson estimates
are uniform.  In particular, on a fixed-radius ball in the rescaled
spaces below there is a number \(d_*>1\) and a constant \(C\) such
that
\begin{equation}\label{co:dirichlet-sobolev}
 \|h\|_{L^2}+\|h\|_{L^{2d_*}}
 \le C\|\nabla h\|_{L^2},\qquad h\in W^{1,2}_0(B),
\end{equation}
whenever the ball and a fixed larger ball are contained in the
working region.
\end{lemma}

\begin{proof}
The full curvature-dimension condition implies the reduced condition
used in some of the references, by \cite[Proposition~2.5(i)]{BS10}.
The interior Poisson regularity and gradient estimates apply to a
constant right hand side; see \cite{Kell13,Jiang14}.  More precisely,
if \(8B_R(z)\) is compactly contained in the domain, the quantitative
estimate in \cite[Theorem~3.1]{Jiang14} has the form
\[
 \|\nabla w\|_{L^\infty(B_R(z))}
 \le C\bigl(R^{-1}+\sqrt{c_\kappa(R^2)}\bigr)
       \bigl(\|w\|_{L^\infty(B_{8R}(z))}
                  +R^2\|\Delta w\|_{L^\infty(B_{8R}(z))}\bigr).
\]
Here \(c_\kappa\) is the heat-semigroup variance constant in that
estimate; the RCD curvature bound controls it, and the doubling and
Poincare constants in \(C\) are uniform on the rescaled balls in
use.  With bounded rescaled curvature this gives the usual
\(R^{-1}\|w\|_\infty+R|a|\) bound.  The interior local-boundedness
estimate first supplies the \(L^\infty\) norm from an \(L^2\) bound.
Our later uses always leave the required eightfold interior room.
The local test
class has the characterization
\[
 \operatorname{Test}_{\rm loc}(Y)
 =\{f\in D_{\rm loc}(\Delta)\cap L^\infty_{\rm loc}:
       |\nabla f|\in L^\infty_{\rm loc},\ 
       \Delta f\in W^{1,2}_{\rm loc}\},
\]
with the analogous definition on \(\Omega\); see
\cite[Lemmas~2.6--2.7]{CDNZPSW}.  Thus constant-Laplacian solutions
satisfy the stated local test condition.  Equivalently, one may
first multiply by a cutoff to obtain Laplacian-domain and Hessian
regularity, and then use a test cutoff on a smaller ball to obtain
the local product calculus.  The relevant Hessian and product rules
are those of \cite{Gigli18}.

The noncollapsed hypothesis implies that the tangent module has
dimension \(n\) almost everywhere, and that
\(\tr\Hess w=\Delta w\) there.  The dimensional Bochner inequality
therefore implies \(M(w)\ge|\Hess w|^2\m\); see
\cite{DPG18,Han18}.  Expanding the square
\[
 |\Hess w-bg|^2=|\Hess w|^2-2b\Delta w+nb^2
\]
proves \eqref{co:shifted-bochner}.  Local finiteness and the
measure formulation follow from the same local test calculus.
Polarization of a nonnegative quadratic form gives the last
assertion about \((w,b)\).

The cutoff, Poincare, and Sobolev statements are local RCD
estimates.  For the heat-regularized cutoff construction, see
\cite[Lemma~3.1]{MN19} and \cite[Lemma~2.6]{CDNZPSW}.
To obtain the zero-boundary coercivity in
\eqref{co:dirichlet-sobolev}, extend \(h\) by zero to a fixed larger
ball.  Poincare on that ball and the positive lower bound for the
measure of its annulus give the \(L^2\) estimate.  The local Sobolev
inequality followed by this estimate gives the other term.  One
may use an effective doubling dimension larger than both \(n\)
and \(2\), so that \(d_*>1\) also when \(n=2\); see
\cite[Theorem~5.1 and the following Remark~3]{HK00}. The annulus lower
bound and all constants are uniform for the almost Euclidean balls
used below.
\end{proof}

For a column \(u=(u_1,\ldots,u_n)^T\) of harmonic functions, define
\begin{equation}\label{co:G-M-matrix}
 G=(\Gamma(u_j,u_k))_{j,k=1}^n,
 \qquad M=(M(u_j,u_k))_{j,k=1}^n,
\end{equation}
where \(M(\cdot,\cdot)\) denotes polarization of
\eqref{co:M-definition}.  Thus \(M[b]=M(b\cdot u)\) is a positive
measure for every \(b\in\R^n\).  If \(E\) is Borel, positivity
implies
\begin{equation}\label{co:matrix-measure-cs}
 |M(v,w)|(E)\le \sqrt{M[v](E)M[w](E)}.
\end{equation}
Here \(|M(v,w)|\) is total variation.  To see this, take the trace
measure as a dominating measure and apply matrix Cauchy--Schwarz
to the positive semidefinite matrix density.  The same argument
applies to integrals with a nonnegative weight.

We record precisely the rescaling regime used throughout the proof.
\(\omega_n\) denotes the measure of the unit ball in \(\R^n\).

\begin{lemma}[An expanding almost Euclidean regime]\label{co:flat-setup}
Fix an \(n\)-regular point \(p\) in the space of
Theorem~\ref{thm:main}.  There are rescalings
\[
 (X_\nu,d_\nu,\m_\nu,p_\nu)
 =(X,s_\nu^{-1}d,s_\nu^{-n}\mathcal H_d^n,p),
 \qquad s_\nu\downarrow0,
\]
numbers \(\kappa_\nu\downarrow0\), and radii \(T_\nu\uparrow\infty\)
with \(\kappa_\nu T_\nu^2\to0\), having the following properties.
The space \(X_\nu\) is \(\RCD(-\kappa_\nu,n)\).  For a sequence
\(\delta_\nu\downarrow0\), all balls with center in
\(B_{64T_\nu}(p_\nu)\) and radius at most \(64T_\nu\) satisfy
\begin{equation}\label{co:uniform-flatness}
 \left|\frac{\m_\nu(B_r(y))}{\omega_n r^n}-1\right|
 \le\delta_\nu,
 \qquad
 d_{GH}\bigl(\overline B_r(y),\overline B_r(0^n)\bigr)
 \le\delta_\nu r.
\end{equation}
There are harmonic columns \(u_\nu\) on
\(B_{32T_\nu}(p_\nu)\), converging on every fixed ball locally
uniformly and strongly in \(W^{1,2}\) to Euclidean coordinates, and
\begin{equation}\label{co:original-gradient-bound}
 |\nabla u_\nu|\le C_n
 \quad\text{on }B_{16T_\nu}(p_\nu).
\end{equation}
\end{lemma}

\begin{proof}
Take rescalings tending to zero.  Compactness and the regular-point
hypothesis imply pointed convergence to \(\R^n\), after passage to
a subsequence, and noncollapsed volume convergence identifies the
limit measure with \(\mathcal H^n\); see \cite{DPG18}.  The rescaled
curvature lower bound can be taken with
\(\kappa_\nu=\max\{-K,0\}s_\nu^2\), enlarged by an arbitrarily
small positive number if necessary.

Fix for the moment a large finite radius.  At every bounded center,
volumes of balls at a fixed positive radius converge to their
Euclidean values.  This convergence is uniform in the center:
otherwise a subsequence of centers converges in the Euclidean
limit and contradicts volume convergence at that limit center.
For a noncollapsed RCD space the volume density is at most one at
every point, not merely almost everywhere.  Bishop--Gromov
therefore traps every smaller model-normalized volume ratio
between its almost maximal value at the fixed radius and one.
The model-volume correction tends uniformly to zero because the
rescaled curvature does.  This proves the first assertion in
\eqref{co:uniform-flatness}, uniformly over all smaller radii.
The almost-volume-rigidity theorem in \cite{DPG18}, applied on a
fixed enlargement of the ball, proves the second assertion.  Its
interior-radius conclusion is why an enlargement is taken first.
No openness of the set of exactly regular points is used.

Interior harmonic approximation to Euclidean coordinate functions
is available by \cite[Corollary~4.12]{AH18}.  For each fixed large
radius construct the harmonic functions on a strictly larger
ball.  Interior energy convergence, local boundedness, and the
interior gradient estimate give convergence to coordinates and a
dimensional gradient bound on the smaller ball.  The latter bound
is independent of the fixed large radius: the coordinate functions
have size \(O(r)\) on a ball of radius \(r\), and the interior
gradient estimate includes the factor \(r^{-1}\).

Finally use a diagonal subsequence, allowing the finite radii just
chosen to increase sufficiently slowly.  It can be chosen so that
all the estimates hold up to the multiples of \(T_\nu\) displayed
above, with \(\delta_\nu\to0\) and
\(\kappa_\nu T_\nu^2\to0\).  This gives
\eqref{co:original-gradient-bound} and the expanding domains of
harmonicity simultaneously.
\end{proof}

From now on the sequence index is suppressed.  Every assertion on
a fixed ball, or with finitely many fixed enlargement factors, is
understood to hold for all sufficiently large \(\nu\).  A quantity
denoted by \(o(1)\) is uniform in the centers in that ball and in
all the smaller scales.  We set
\begin{equation}\label{co:dyadic-R}
 r_i=2^{-i},\qquad
 R_i(y)=\fint_{B_{r_i}(y)}G\dd\m,
 \qquad i=0,1,\ldots.
\end{equation}

\subsection{Transformation at every center and scale}

Harmonic normalization across scales is a central tool in
\cite[Theorem~1.32]{CN15}, \cite[Section~7]{CJN21}, and
\cite[Theorem~4.1]{HH24}. The noncollapsed RCD transformation framework of
Bru\`e--Naber--Semola \cite[Section~3.3, Proposition~3.13 and Corollary~3.16]{BNS20}
is a further methodological antecedent. We give the all-center statement and
its proof needed on the expanding domains used here.

\begin{lemma}[Uniform harmonic normalization]\label{co:harmonic-normalization}
For every pair of fixed constants \(A,C\ge1\), the matrices \(R_i(y)\) are
positive definite for \(y\in B_A(p)\) and \(i\ge0\).  In the
rescaled ball \(B_{Cr_i}(y)\), the map
\begin{equation}\label{co:normalized-column}
 U_{i,y}=\frac{R_i(y)^{-1/2}(u-u(y))}{r_i}
\end{equation}
is, after an orthogonal change of target coordinates, an
\(o(1)\) approximation to Euclidean coordinates, uniformly and
strongly in interior \(W^{1,2}\).  Its gradients are uniformly
bounded on each fixed interior enlargement.  Its Gram matrix
converges strongly in every finite \(L^p\) on such enlargements
to \(\Id\).

For every fixed \(C\ge1\), if \(d(y,z)\le Cr_i\) and
\(C^{-1}r_i\le r_j\le Cr_i\), then
\begin{equation}\label{co:nearby-R}
 (1-o(1))R_i(y)\le R_j(z)\le(1+o(1))R_i(y).
\end{equation}
For every \(a>0\), after choosing a sufficiently late rescaling,
\begin{equation}\label{co:slow-R}
 C_a^{-1}2^{-2a|i-j|}R_j(y)
 \le R_i(y)\le C_a2^{2a|i-j|}R_j(y).
\end{equation}
The constants are uniform in the center and the two indices.  In
particular, since \(R_0(y)=\Id+o(1)\) on fixed balls,
\begin{equation}\label{co:R-lower}
 c_a r_i^{2a}\Id\le R_i(y)\le C\Id.
\end{equation}
\end{lemma}

\begin{proof}
We give the compactness induction, including the order of the
smallness choices.  Say that a harmonic column is
\(\tau\)-normalizable at \((y,r)\) if there is an invertible
matrix \(T\) for which
\(T(u-u(y))/r\), on \(B_{2r}(y)\) in units \(r\), differs by at
most \(\tau\) from a Euclidean coordinate map under a
Gromov--Hausdorff approximation.  Replacing the uniform error in
this definition by an interior normalized \(L^2\) error gives an
equivalent induction, after changing its tolerance by a fixed
factor; we use harmonic interior estimates on a slightly larger
ball whenever needed.  Rotations of the Euclidean coordinates
are allowed.

Choose a small fixed tolerance \(\tau\).  In an exactly Euclidean
ball, a harmonic column \(V\) that is \(\tau\)-close to the
coordinate functions has
\[
 \|DV(0)-\Id\|\le C_n\tau,
 \qquad
 \sup_{B_{2b}}|DV(0)^{-1}(V-V(0))-x|
 \le C_n\tau b^2.
\]
These follow by applying interior derivative estimates to \(V-x\).
Choose a dyadic number \(b=2^{-m}\) so small that
\(2C_n b<1/4\), and then choose \(\tau\) small enough that the
inverse derivative is bounded.  Thus affine renormalization
improves the error in units \(b\) by a strict factor.

The same propagation holds on all sufficiently flat balls in our
sequence.  Indeed, if it failed, rescale a failing ball to unit
size and apply its assumed normalization.  The spaces converge to
a Euclidean ball by \eqref{co:uniform-flatness}.  The normalized
harmonic columns have uniform local \(L^2\), energy, and gradient
bounds and therefore converge locally uniformly and strongly in
interior \(W^{1,2}\) to a Euclidean harmonic column.  These are
interior convergence statements from \cite[Section~4]{AH18}; no
convergence of a full boundary-value problem is required.  The
limit has the \(\tau\)-normalizability property.  Its inverse
derivative supplies a fixed matrix on the sequence which gives
error strictly less than \(\tau\) at radius \(b\), a
contradiction.  This additional correction to the already normalized
target, and its inverse, have norm at most \(1+C_n\tau\), after
allowing the orthogonal target changes in the definition.

At each bounded center the initial normalizability is supplied by
Lemma~\ref{co:flat-setup}, uniformly in that center.  Run the
preceding step down each of the finitely many dyadic residue
classes modulo \(m\).  This proves normalizability at every
dyadic scale.  The mean Gram matrix of a normalized column on its
unit ball is \(\Id+O(\tau)\), uniformly once the space is
sufficiently flat: a contrary sequence again converges strongly
in energy to a Euclidean harmonic column, whose derivative is
\(O(\tau)\)-close to \(\Id\) on that interior ball.  Consequently
these Gram matrices are nonsingular and replacing an inductively
chosen \(T\) by the canonical normalization \(R_i(y)^{-1/2}\)
changes the normalized map only by an orthogonal map and an
\(O(\tau)\) error.

To see explicitly why fixed enlargements are permitted, fix their
largest factor first.  Start the preceding induction on a larger
fixed radius, allowing finitely many negative dyadic indices.
For any scale \(r_i\), an ancestor radius contains the desired
enlargement, and at most a fixed number of the contraction steps
separate that ancestor from \(r_i\).  Their normalizing matrices
and inverse matrices are uniformly comparable.  Apply the
interior estimates on the ancestor ball and identify the mean
Gram normalization on the descendant ball by strong energy
convergence.  All factors and the number of steps are fixed before
\(\tau\) and the geometric flatness are chosen.  This proves
normalizability and bounded normalized gradients on every fixed
enlargement.  As \(\tau\) is arbitrary, the errors tend uniformly
to zero along the original sequence, including when the chosen
center and scale depend on the sequence index.

The same argument on an enlargement containing both balls in
\eqref{co:nearby-R} proves that their Gram averages in the first
normalization both tend to \(\Id\).  This gives
\eqref{co:nearby-R}.  Iterating its same-center adjacent-scale
case, with error chosen so that each factor is at most \(2^{2a}\),
gives \eqref{co:slow-R}.  The upper bound in
\eqref{co:R-lower} also follows directly from
\eqref{co:original-gradient-bound}.

For completeness, the strong convergence of the Gram functions
can be seen without assuming convergence of pointwise gradient
representatives.  A Bochner cutoff estimate first bounds the
local normalized Hessian energy.  The Hessian product rule and
bounded normalized gradients then bound the Gram functions in
local \(W^{1,2}\).  Sobolev compactness and the convergence of
integrals of gradient scalar products identify their limit with
\(\Id\); see \cite[Corollary~4.3 and Theorem~4.6]{AH18}.  Their
uniform boundedness upgrades the resulting strong \(L^2\)
convergence to every finite \(L^p\).
\end{proof}

\begin{lemma}[Small normalized Bochner mass]\label{co:bochner-small}
For each pair of fixed constants \(A,C\ge1\), uniformly for \(y\in B_A(p)\) and
\(i\ge0\),
\begin{equation}\label{co:small-mass}
 r_i^{2-n}M(B_{Cr_i}(y))\le o(1)R_i(y).
\end{equation}
\end{lemma}

\begin{proof}
In units \(r_i\), apply Lemma~\ref{co:local-calculus} to a scalar
combination of \(U=U_{i,y}\).  The rescaled Bochner matrix is
\begin{equation}\label{co:M-scaling}
 M_{r_i}(U)=r_i^{2-n}R_i(y)^{-1/2}M(u)R_i(y)^{-1/2},
\end{equation}
where the curvature term in its definition is
\(\kappa r_i^2\).  Choose a nonnegative cutoff equal to one on
\(B_C\) and supported on a larger fixed ball.  Its gradient and
Laplacian are uniformly bounded.  For a fixed vector \(b\),
\[
 \int\phi\dd M_{r_i}(b\cdot U)
 =\frac12\int\bigl(\Gamma(b\cdot U)-|b|^2\bigr)
                  \Delta\phi\dd\m_{r_i}
  +\kappa r_i^2\int\phi\Gamma(b\cdot U)\dd\m_{r_i}.
\]
Here \(\int\Delta\phi\dd\m_{r_i}=0\).  Both terms tend uniformly
to zero by Lemma~\ref{co:harmonic-normalization}.  Positivity of
the Bochner measure bounds its mass on \(B_C\).  Polarization in
the fixed finite dimension, or the same argument for unit
vectors, gives the matrix bound.  Conjugating back proves
\eqref{co:small-mass}.
\end{proof}

\subsection{Open coordinate images}

\begin{lemma}[Open harmonic coordinates]\label{co:open-coordinates}
For every fixed \(A\ge1\), the map \(u\) is a homeomorphism
from \(B_A(p)\) onto an open subset of \(\R^n\), once the
rescaling is sufficiently late.  There is a fixed \(c>0\) such
that, whenever the balls involved stay in a fixed interior
working region,
\begin{equation}\label{co:ellipsoid-containment}
 \{P:|P-u(x)|_{R_i(x)^{-1}}<c r_i\}
 \subset u(B_{r_i/2}(x)).
\end{equation}
Moreover, for every fixed \(D\ge1\), if
\(D^{-1}r_i\le d(x,z)\le Dr_i\), then
\begin{equation}\label{co:separation-scale}
 \bigl||u(z)-u(x)|_{R_i(x)^{-1}}-d(x,z)\bigr|
 \le o_D(1)r_i.
\end{equation}
\end{lemma}

\begin{proof}
The approximation in Lemma~\ref{co:harmonic-normalization} gives
\eqref{co:separation-scale}.  Test a pair of distinct points at a
dyadic scale comparable with its distance.  Once the error is
less than one half of that distance, the coordinate values are
distinct.  Distances larger than the initial unit scale in a
fixed ball are separated directly by the initial uniform
Euclidean approximation.  Thus \(u\) is injective on \(B_A(p)\).

We prove \eqref{co:ellipsoid-containment} with its interior margin.
Fix a target \(P\) in the indicated ellipsoid and put
\(x_0=x\), \(\rho_0=r_i\).  Write \(R(\rho,y)\) for the Gram
average at a dyadic radius.  Suppose inductively that
\[
 |P-u(x_j)|_{R(\rho_j,x_j)^{-1}}<c\rho_j,
 \qquad \rho_j=4^{-j}\rho_0.
\]
Approximate surjectivity of the normalized Euclidean coordinate
map provides \(x_{j+1}\) with
\[
 d(x_j,x_{j+1})\le(c+\eta)\rho_j,
 \qquad
 |P-u(x_{j+1})|_{R(\rho_j,x_j)^{-1}}\le\eta\rho_j.
\]
Here \(\eta\to0\) is uniform in the scale.  For example, choose a
point under the Gromov--Hausdorff correspondence whose Euclidean
coordinate is the normalized target.  Its distance from the
center and the error in its coordinate value have the displayed
bounds.  By \eqref{co:nearby-R}, the second bound is at most
\(2\eta\rho_j\) when measured in
\(R(\rho_j/4,x_{j+1})^{-1}\).  Choose \(c<1/8\), and then choose
\(\eta<c/8\).  It is at most \(c\rho_{j+1}\), so the induction
continues.

The total displacement is at most
\((c+\eta)\sum_j\rho_j<\rho_0/2\).  All centers therefore stay
in the region where the estimates are valid.  Completeness gives
a limit \(x_\infty\in B_{\rho_0/2}(x)\), with a strict margin
because of the preceding choices.  The upper bound \(R(\rho,y)
\le C\Id\) in the original coordinate units shows that the
Euclidean residual \(|P-u(x_j)|\) tends to zero.  Continuity of
\(u\) now gives \(u(x_\infty)=P\).

Apply the ellipsoid statement at each point with a sufficiently
small radius contained in \(B_A(p)\).  Its image is open.  Finally,
a continuous injection is a homeomorphism locally on each compact
interior ball, since these spaces are proper.  Alternatively the
same separation-scale estimate gives continuity of the inverse.
This proves the assertion on the whole open domain.
\end{proof}

\section{Heat averages and cumulative matrix control}
\label{co:heat-section}

We next replace the ball Gram averages by matrices that are
monotone as the spatial scale decreases.  The local mass forms use
inverse ball averages at nearby centers.  Relative comparability
of these averages with the heat averages does not by itself compare
the twice-conjugated mass forms.
We will bound the accumulated cost of changing the normalization by
a small multiple of the final inverse heat matrix.
Heat Gram normalization and quadratic telescoping are used in
\cite[Sections~8.1--8.2]{CJN21}. The estimates below give the cumulative
comparison for the inverse matrices without a commutativity
assumption.

\subsection{Continuous kernels and compactly supported sources}

Let \(p_t(x,z)\) be the common continuous heat kernel and \(P_t\)
the heat semigroup.  We use the Gaussian and first spatial
gradient estimates of \cite[Theorem~1.2 and Corollary~1.2]{JLZ16}.
On the fixed testing regions, for \(0<t\le1\), volume comparison
allows their form
\begin{equation}\label{co:kernel-bounds}
 p_t(x,z)\le C t^{-n/2}e^{-d(x,z)^2/(Ct)},
 \qquad
 |\nabla_xp_t(x,z)|\le C t^{-(n+1)/2}e^{-d(x,z)^2/(Ct)}.
\end{equation}
The gradient statement initially holds almost everywhere.  The
lower bound, used when \(d(x,z)\le A\sqrt t\), is
\begin{equation}\label{co:kernel-lower}
 p_t(x,z)\ge c_A t^{-n/2}.
\end{equation}
We may use one fixed weaker negative curvature bound in these
estimates.  The constants are then uniform in the sequence.

The following interpretation will be useful because \(M\) can
have a singular part.

\begin{lemma}[Kernel differences and measure sources]\label{co:kernel-measures}
Fix \(C_0\ge1\), \(x,y\) in a fixed testing region, and
\(d(x,y)\le C_0r\), where \(0<r\le1\).  Put
\[
 K_t(z)=\fint_{B_r(y)}p_t(a,z)\dd\m(a),
 \qquad D=d(x,z).
\]
After enlarging fixed constants, the following bounds hold at
every source pole \(z\):
\begin{align}
 K_t(z)&\le C r^{-n}e^{-cD^2/r^2},&&0<t\le r^2,
       \label{co:short-kernel}\\
 |K_t(z)-p_t(x,z)|&\le C r t^{-(n+1)/2}e^{-cD^2/t},
       &&r^2\le t\le1.\label{co:long-kernel}
\end{align}
If \(f\in W^{1,2}\cap L^\infty\) has compact support and its
Laplacian is a finite signed measure \(\mu\) of compact support,
then
\begin{equation}\label{co:semigroup-measure}
 P_t f-P_s f=\int_s^t P_\tau\mu\dd\tau,
 \qquad 0<s<t,
\end{equation}
using the continuous representatives.  In the application
\(f=\chi\Gamma(u_j,u_k)\), this identity can also be averaged over
a ball and passed to \(s\downarrow0\).
\end{lemma}

\begin{proof}
When \(D\le 2(C_0+2)r\), symmetry and total heat mass bound the
integral of \(p_t(a,z)\) over the ball by one.  Dividing by its
volume gives \(Cr^{-n}\), which is the first bound with adjusted
constants.  If \(D>2(C_0+2)r\), then \(d(a,z)\ge D/2\) for all
\(a\in B_r(y)\).  The Gaussian upper bound gives
\(Ct^{-n/2}e^{-cD^2/t}\).  For \(t\le r^2\) and \(D/r\) bounded
below by this fixed constant, its small-time power is absorbed
in the exponential, giving \eqref{co:short-kernel}.

For the second bound, \(d(a,x)\le(C_0+1)r\).  Integrate the first
spatial gradient estimate along a minimizing curve from \(x\) to
\(a\); that curve stays in a fixed testing enlargement.  The
inequality
\(d(b,z)^2\ge D^2/2-d(b,x)^2\), with \(t\ge r^2\), absorbs the
variation of the Gaussian along the curve.  This gives
\eqref{co:long-kernel}, initially for almost every fixed source
pole \(z\).  For each such pole, local Sobolev-to-Lipschitz gives
the difference inequality for the continuous representative in
the first variable.  By symmetry, continuity of the common heat
kernel, and full support, a sequence of these full-measure poles
converges to any prescribed \(z\).  Passing to the limit, with
relaxed Gaussian constants if necessary, proves the bounds at
every pole.  No gradient representative has been paired with a
singular measure.

For \eqref{co:semigroup-measure}, first test the distributional
Laplacian against positive-time heat regularizations of compactly
supported test functions.  Symmetry and integration by parts give
the asserted identity in distributions.  Equivalently, differentiate
\(P_\tau f\) for positive \(\tau\), pair the Laplacian with the
continuous Sobolev function \(p_\tau(x,\cdot)\), and integrate.
A cutoff equal to one on the compact supports justifies this
pairing; the kernel belongs locally to \(W^{1,2}\) at positive
time.  The measure and its Sobolev-dual realization agree on this
continuous representative.  The kernel bounds give continuity of
both sides in \(x\), so the identity holds pointwise.

In our application the short-time averaged integral of the total
variation of the source is finite: \eqref{co:short-kernel},
integrated up to \(r^2\), bounds it by a finite Gaussian-weighted
multiple of \(r^{2-n}|\mu|\).  Also \(P_s f\to f\) in \(L^1\)
on the averaging ball.  Thus the averaged identity passes to
\(s=0\) by absolute convergence.  These statements also follow
by first using positive lower time endpoints throughout and then
applying the bounds just proved.
\end{proof}

Choose a nonnegative cutoff \(\chi\) equal to one on \(B_T(p)\),
supported in \(B_{2T}(p)\), with
\[
 |\nabla\chi|\le C/T,\qquad |\Delta\chi|\le C/T^2,
\]
where \(T=T_\nu\) comes from Lemma~\ref{co:flat-setup}.  Extend
\(\chi G\) by zero outside its support.  Its distributional
Laplacian is the compactly supported matrix measure
\begin{equation}\label{co:cutoff-source}
 \Delta(\chi G)=2\chi M-2\kappa\chi G\m+E_{\rm cut}\m,
 \qquad E_{\rm cut}=(\Delta\chi)G+2\Gamma(\chi,G).
\end{equation}
The last expression is a matrix of \(L^1\) functions supported in
the cutoff annulus.

\begin{lemma}[Monotone heat Gram matrices]\label{co:heat-matrices}
There is \(\lambda=\lambda_\nu>0\), with \(\lambda_\nu\to0\),
such that, at all fixed testing centers and \(0<t\le1\),
\begin{equation}\label{co:heat-derivative}
 P_t(\chi M)
 \le\partial_t\bigl(P_t(\chi G)+\lambda t\Id\bigr)
 \le 3P_t(\chi M)+C\Id.
\end{equation}
Define
\begin{equation}\label{co:QH-definition}
 Q_i(x)=P_{r_i^2}(\chi G)(x)+\lambda r_i^2\Id,
 \qquad H_i(x)=Q_i(x)^{-1}.
\end{equation}
Then
\begin{align}
 &(1-o(1))R_i(x)\le Q_i(x)\le(1+o(1))R_i(x),
       \qquad Q_0(x)=\Id+o(1),\label{co:QR-comparison}\\
 &H_i(x)\le H_{i+1}(x)\le(1+o(1))H_i(x).
       \label{co:H-monotone}
\end{align}
The nearby-center and slow-growth estimates of
Lemma~\ref{co:harmonic-normalization} hold for \(Q_i\) as well.
In particular, for every prescribed \(a>0\), with a sufficiently
late rescaling,
\begin{equation}\label{co:slow-QH}
 Q_j(x)\le C_a2^{2a(i-j)}Q_i(x)\quad(j\le i),
 \qquad \|H_i(x)\|\le C_a2^{2ai}.
\end{equation}
\end{lemma}

\begin{proof}
The original gradient bound and a Bochner cutoff on a ball of
radius \(s\ge1\) give, for centers in a fixed bounded set and up
to the support region,
\begin{equation}\label{co:raw-mass-growth}
 M(B_s(x))[v]\le C(s^{n-2}+\kappa s^n)|v|^2
 \le Cs^{n-2}|v|^2.
\end{equation}
For radii larger than the support radius the same bound holds
with \(\chi M\) in place of \(M\), by monotonicity and \(n\ge2\).
The cutoff for this estimate lies inside the expanding domain
where \eqref{co:original-gradient-bound} holds.  The same estimate
for the Hessian, followed by Cauchy--Schwarz in the product rule,
gives
\[
 \int|E_{\rm cut}|\dd\m\le CT^{n-2}.
\]
Since its support is a distance at least \(T/2\) from the fixed
testing centers for large \(\nu\), \eqref{co:kernel-bounds} yields
\begin{equation}\label{co:cutoff-small}
 \sup_{0<t\le1}\|P_t E_{\rm cut}(x)\|
 \le C T^{n-2}\sup_{0<t\le1}t^{-n/2}e^{-cT^2/t}
 =o(1).
\end{equation}
Choose \(\lambda\) at least this supremum plus
\(2\kappa\|G\|_\infty\), and add a positive number tending to
zero.  Equations \eqref{co:semigroup-measure} and
\eqref{co:cutoff-source} give
\[
 \partial_t(P_t(\chi G)+\lambda t\Id)
 =2P_t(\chi M)-2\kappa P_t(\chi G)
     +P_tE_{\rm cut}+\lambda\Id.
\]
The last three terms have nonnegative sum and bounded norm.  This
proves \eqref{co:heat-derivative}, in fact with lower bound
\(2P_t(\chi M)\).

We verify relative, rather than merely absolute, comparison with
the ball averages.  Fix \(i,x\) and conjugate by \(R_i(x)^{-1/2}\).
On every fixed ball \(B_{Lr_i}(x)\), the conjugated Gram field
converges in normalized \(L^1\) to \(\Id\), by
Lemma~\ref{co:harmonic-normalization}.  The heat kernel at time
\(r_i^2\) is at most \(Cr_i^{-n}\), so the integral of its local
error tends to zero.  The tails are uniform in \(i\): on the
shell with radius \(2^m r_i\le1\), slow growth bounds the
conjugated Gram average by \(C2^{2am}\), while the Gaussian
bounds its contribution by
\[
 C2^{(n+2a)m}e^{-c4^m}.
\]
When \(2^m r_i>1\), use the original gradient bound and
\eqref{co:R-lower}; its conjugation cost \(Cr_i^{-2a}\) is at
most \(C2^{2am}\).  Volume comparison and the same Gaussian again
give the displayed summable bound.  This treats the entire
support of \(\chi\).  The missing mass from \(1-\chi\) is
bounded in the same way, using total heat mass one.  First choose
\(L\) large enough to make the tail small and then choose the
rescaling late enough to make the local error small.  Beyond the
expanding almost Euclidean region, the global volume bound adds
at most \(e^{C2^m r_i}\) to a shell estimate; since \(r_i\le1\),
this factor is absorbed by decreasing the Gaussian constant in
\(e^{-c4^m}\).  Thus the assertion about the missing heat mass is
uniform there as well.  Finally,
\[
 \lambda r_i^2R_i(x)^{-1}\le C_a\lambda r_i^{2-2a}\Id=o(1)\Id
\]
for \(a<1\).  This proves \eqref{co:QR-comparison}; its initial
scale assertion follows from Euclidean convergence.

The derivative positivity makes \(Q_i\ge Q_{i+1}\), hence
\(H_i\le H_{i+1}\).  Relative comparison with \(R_i\) and
\eqref{co:nearby-R} give the adjacent-scale upper bound and all
fixed nearby comparisons.  Iteration gives
\eqref{co:slow-QH}, choosing the adjacent-scale tolerance after
the prescribed exponent \(a\).
\end{proof}

\subsection{The dyadic heat comparison}

The matrices
\begin{equation}\label{co:N-definition}
 N_j(x)=r_j^{2-n}M(B_{r_j}(x))
\end{equation}
are masses in the original harmonic coordinate components.  They
have not been conjugated by a Gram inverse.  In particular,
Lemma~\ref{co:bochner-small} and \eqref{co:QR-comparison} give
\begin{equation}\label{co:N-small}
 N_j(x)\le\varepsilon_\nu Q_j(x),
 \qquad \varepsilon_\nu\longrightarrow0,
\end{equation}
uniformly in \(j\) and in fixed testing regions.  We allow
\(\varepsilon_\nu\) to increase by fixed multiplicative constants
in what follows.

\begin{lemma}[Dyadic heat comparison]\label{co:heat-comparison}
Fix constants \(A,C_0\ge1\), where \(C_0\) bounds
\(d(x,y)/r_i\), and fix a bounded region for \(x\).
There is a fixed exponent \(\gamma>0\), which
may be taken to be \(1/2\), with the following properties for all
sufficiently late rescalings.  Set
\begin{equation}\label{co:F-definition}
 \mathcal F_i(x)=2^{-\gamma i}\Id+
       \sum_{j=0}^i2^{-\gamma(i-j)}N_j(x).
\end{equation}
For the prescribed enlargement \(A\),
\begin{equation}\label{co:heat-decrement}
 c_A r_i^{2-n}M(B_{Ar_i}(x))
 \le Q_i(x)-Q_{i+1}(x)\le C\mathcal F_i(x).
\end{equation}
For \(d(x,y)\le C_0r_i\) and \(v,w\in\R^n\),
\begin{equation}\label{co:bilinear-heat-comparison}
 |v^T(R_i(y)-Q_i(x))w|
 \le C|v|_{Q_i(x)}\sqrt{\mathcal F_i(x)[w]}.
\end{equation}
All constants are independent of \(i\) and of the expanding
cutoff radius.  They may depend on the stated fixed parameters.
\end{lemma}

\begin{proof}
\emph{The semigroup identity.}
Write \(r=r_i\), suppress \(x\) in \(N_j,Q_j\), and let
\(\mathcal A f=\fint_{B_r(y)}f\dd\m\).  We first prove the more
explicit estimate
\begin{equation}\label{co:raw-bilinear-heat}
 \begin{split}
 |v^T(R_i(y)-Q_i(x))w|
 &\le Cr|v||w|\\
 &\quad+C\sum_{j=0}^i2^{-(i-j)}
                   \sqrt{N_j[v]N_j[w]}.
 \end{split}
\end{equation}
This also specifies all the spatial and time tails entering
\eqref{co:bilinear-heat-comparison}.

Put \(f=\chi G\), \(S=\Delta f\), and use the kernel functions
from Lemma~\ref{co:kernel-measures}.  Since \(\chi=1\) on the
averaging ball, the exact semigroup identity is
\begin{equation}\label{co:average-heat-identity}
 \begin{split}
 R_i(y)-Q_i(x)
 ={}&\mathcal A P_1f-P_1f(x)-\lambda r^2\Id\\
 &-\int_0^{r^2}\mathcal A P_tS\dd t
  -\int_{r^2}^1\bigl(\mathcal A P_tS-P_tS(x)\bigr)\dd t.
 \end{split}
\end{equation}
The averaged endpoint at zero is justified in
Lemma~\ref{co:kernel-measures}; it does not require pointwise
convergence of the Gram field.

\smallskip\noindent
\emph{The local source.}
Consider the contribution of \(2\chi M\) on \(B_1(x)\).
The integrated short-time kernel obeys
\begin{equation}\label{co:integrated-short-kernel}
 \int_0^{r^2}K_t(z)\dd t
 \le Cr^{2-n}e^{-c d(x,z)^2/r^2}.
\end{equation}
For \(0\le q<i\), integration of
\eqref{co:long-kernel} on the time block
\(I_q=[r_{q+1}^2,r_q^2]\) gives
\begin{equation}\label{co:integrated-long-kernel}
 \int_{I_q}|K_t(z)-p_t(x,z)|\dd t
 \le Cr r_q^{1-n}e^{-c d(x,z)^2/r_q^2}.
\end{equation}
To integrate a Gaussian at scale \(r_q\) against \(M(v,w)\),
partition \(B_1(x)\) into \(B_{r_q}(x)\) and the annuli
\[
 B_{r_j}(x)\setminus B_{r_{j+1}}(x),\qquad 0\le j<q.
\]
Define positive numerical coefficients
\begin{equation}\label{co:gaussian-shell-coefficients}
 k_0=1,\qquad
 k_h=2^{(n-2)h}e^{-c4^{h-1}}\quad(h\ge1).
 \qquad \sum_{h=0}^\infty2^h k_h<\infty.
\end{equation}
On the annulus indexed by \(j<q\), the distance from \(x\) is at
least \(2^{q-j-1}r_q\).  Equation
\eqref{co:matrix-measure-cs} and the identity
\(M(B_{r_j}(x))=r_j^{n-2}N_j\) show that the contribution of the
time block \(I_q\) is at most
\[
 C2^{-(i-q)}\sum_{j=0}^q k_{q-j}
                       \sqrt{N_j[v]N_j[w]}.
\]
Reversing the finite sums, the sum of all these contributions is
at most
\begin{align*}
 &C\sum_{j=0}^{i-1}2^{-(i-j)}
       \left(\sum_{h=0}^{i-1-j}2^h k_h\right)
                       \sqrt{N_j[v]N_j[w]}\\
 &\hspace{25mm}\le
 C\sum_{j=0}^{i-1}2^{-(i-j)}\sqrt{N_j[v]N_j[w]}.
\end{align*}
The short-time term uses the same partition with \(q=i\) and
\eqref{co:integrated-short-kernel}.  It is bounded by
\[
 C\sum_{j=0}^i k_{i-j}\sqrt{N_j[v]N_j[w]}
 \le C\sum_{j=0}^i2^{-(i-j)}\sqrt{N_j[v]N_j[w]}.
\]
This proves the required contribution on \(B_1(x)\), with the
original unit-ball \(N_j\) and no unmentioned enlarged mass.

\smallskip\noindent
\emph{The remaining source terms.}
We treat all remaining parts of
\eqref{co:average-heat-identity}.  If \(D=d(x,z)\ge1\), the
integrated kernel bounds imply
\begin{equation}\label{co:exterior-integrated-kernel}
 \int_0^{r^2}K_t(z)\dd t+
 \int_{r^2}^1|K_t(z)-p_t(x,z)|\dd t
 \le Cr e^{-c'D^2}.
\end{equation}
For the second term, split the Gaussian into two factors, use
one to retain \(e^{-c'D^2}\), and use the other to integrate
\(t^{-(n+1)/2}e^{-c''/t}\) on \((0,1)\).  For the first term,
\eqref{co:integrated-short-kernel} reduces the assertion to
absorbing \(r^{1-n}\) in \(e^{-c''/r^2}\).  Both bounds are uniform
for \(0<r\le1\).

Equation \eqref{co:raw-mass-growth}, summed over the annuli
\(B_{2^{h+1}}(x)\setminus B_{2^h}(x)\), gives
\[
 \int e^{-c'd(x,z)^2}\dd(\chi M[v])(z)\le C|v|^2.
\]
Therefore \eqref{co:exterior-integrated-kernel} and polarized
Cauchy--Schwarz bound the exterior positive-source contribution
by \(Cr|v||w|\).  The total variation of the cutoff source is
\(O(T^{n-2})\), is supported a distance comparable with \(T\)
from \(x\), and has the same Gaussian moment bound, in fact a
bound tending to zero.  Its contribution is also at most
\(Cr|v||w|\).

The curvature source in \eqref{co:cutoff-source} has bounded
density, of norm at most \(C\kappa\).  For its short-time part use
heat mass one.  For its long-time part, integration of
\eqref{co:long-kernel} in the source variable gives
\(Cr/\sqrt t\).  This last Gaussian integral follows from local
volume comparison for nearby shells and the global exponential
volume bound for remote shells.  Consequently the curvature
contribution is at most
\[
 C\kappa\left(r^2+\int_{r^2}^1 r t^{-1/2}\dd t\right)|v||w|
 \le Cr|v||w|.
\]
The time-one heat regularization of the bounded matrix field
\(f\) has uniformly bounded spatial gradient in the testing
region, by the integrated first kernel gradient estimate.
Hence \(|v^T(\mathcal A P_1f-P_1f(x))w|\le Cr|v||w|\).
Also \(\lambda r^2|v||w|\le Cr|v||w|\).  These facts prove
\eqref{co:raw-bilinear-heat}.

\smallskip\noindent
\emph{Normalization and decrements.}
Choose any sufficiently small fixed growth exponent \(a>0\) in
\eqref{co:slow-QH}.  By \eqref{co:N-small}, for \(j\le i\),
\[
 \sqrt{N_j[v]}\le C2^{a(i-j)}|v|_{Q_i}.
\]
Weighted Cauchy--Schwarz now gives, for instance when \(a<1/8\),
\begin{align*}
 \sum_{j=0}^i2^{-(i-j)}\sqrt{N_j[v]N_j[w]}
 &\le C|v|_{Q_i}\sum_{j=0}^i2^{-(1-a)(i-j)}\sqrt{N_j[w]}\\
 &\le C|v|_{Q_i}
       \left(\sum_{j=0}^i2^{-(i-j)/2}N_j[w]\right)^{1/2}.
\end{align*}
Indeed the unused scalar series has terms
\(2^{-(2(1-a)-1/2)(i-j)}\), whose sum is bounded.  Since
\(Q_i\ge c_a r_i^{2a}\Id\), the raw term in
\eqref{co:raw-bilinear-heat} is at most
\[
 Cr_i^{1-a}|v|_{Q_i}|w|
 \le C|v|_{Q_i}\bigl(2^{-i/2}|w|^2\bigr)^{1/2}.
\]
This proves \eqref{co:bilinear-heat-comparison} with
\(\gamma=1/2\).

It remains to prove the decrement estimates.  Integrate
\eqref{co:heat-derivative} on
\([r_i^2/4,r_i^2]\).  On a fixed ball \(B_{Ar_i}(x)\), where
\(\chi=1\), the heat lower bound
\eqref{co:kernel-lower} gives the left inequality in
\eqref{co:heat-decrement}.  The upper bound is
\[
 Q_i-Q_{i+1}
 \le Cr_i^{2-n}\int e^{-c d(x,z)^2/r_i^2}\dd(\chi M)(z)
        +Cr_i^2\Id.
\]
On \(B_1(x)\), the shell computation with \(q=i\) bounds the
positive term by \(C\sum_{j=0}^i k_{i-j}N_j\), which is at most
\(C\sum_{j=0}^i2^{-(i-j)/2}N_j\).  Outside \(B_1(x)\),
\eqref{co:raw-mass-growth} and Gaussian decay give \(Cr_i^2\Id\),
after absorbing every fixed power of \(r_i^{-1}\) in
\(e^{-c/r_i^2}\).  Since \(r_i^2\le2^{-i/2}\), this proves the
upper bound with the same \(\gamma\).
\end{proof}

\begin{remark}[Fixed enlargements and indexing]\label{co:enlargement-indexing}
No enlarged mass was needed in the preceding bilinear estimate.
When one is used elsewhere, it can be returned to the same
notation without losing the decay.  Indeed, for fixed \(A\ge1\)
put \(h=\lceil\log_2 A\rceil\).  For \(j\ge h\),
\[
 r_j^{2-n}M(B_{Ar_j}(x))\le2^{h(n-2)}N_{j-h}(x).
\]
The finitely many \(j<h\) are bounded by \(C_A\Id\) in the
original coordinate components, by a fixed-scale cutoff estimate.
Thus an exponentially weighted sum of these enlarged masses is
bounded by a fixed multiple of
\(2^{-\gamma i}\Id+\sum_{j\le i}2^{-\gamma(i-j)}N_j\).
A center shift bounded by \(C_0r_j\) only increases \(A\) by
\(C_0\).  In contrast, the mass matrices \(D_{i,y}\) below include
two Gram inverses; their comparison requires the argument that
follows.
\end{remark}

\subsection{Accumulation without commuting matrices}

We separate the finite-dimensional calculation from its geometric
application.  All inequalities in this subsection are inequalities
of quadratic forms on the fixed space \(\R^n\).

\begin{lemma}[Weighted inverse increments]\label{co:inverse-increments}
Fix a testing center \(x\), and suppress it in the notation.
Let \(a>0\) in \eqref{co:slow-QH} be chosen so that
\(4a<\gamma\).  Set
\[
 B_i=H_{i+1}-H_i,\qquad \Delta_i=Q_i-Q_{i+1},
 \qquad a_i=H_i\xi
\]
for a fixed \(\xi\in\R^n\).  Then, uniformly for every integer
\(k\ge0\),
\begin{align}
 &H_i\Delta_iH_i\le B_i\le(1+C\varepsilon_\nu)H_i\Delta_iH_i,
             \label{co:inverse-increment-order}\\
 &\sum_{i<k}\mathcal F_i[a_i]\le C H_k[\xi],
             \label{co:weighted-total}\\
 &C^{-1}H_k\le \Id+\sum_{i<k}H_iN_iH_i\le C H_k.
             \label{co:central-cumulative}
\end{align}
For every fixed \(A\ge1\), the upper sum bound also holds with
\(r_i^{2-n}M(B_{Ar_i}(x))\) in place of \(N_i\).
\end{lemma}

\begin{proof}
Let \(E_i=Q_i^{-1/2}\Delta_iQ_i^{-1/2}\).
By \eqref{co:H-monotone}, \(0\le E_i\le\varepsilon_\nu\Id\)
after increasing the sequence \(\varepsilon_\nu\to0\).
The exact inverse identity is
\[
 B_i=Q_i^{-1/2}\bigl((\Id-E_i)^{-1}-\Id\bigr)Q_i^{-1/2}.
\]
The spectral inequality
\(E_i\le(\Id-E_i)^{-1}-\Id\le(1+C\varepsilon_\nu)E_i\)
proves \eqref{co:inverse-increment-order}.  This use of the
functional calculus is for one conjugated positive matrix at a
time; matrices from different scales need not commute.

We spell out the convolution estimate.  Since
\(0\le B_l\le C\varepsilon_\nu H_l\), conjugation by
\(Q_l^{1/2}\) gives \(B_lQ_lB_l\le C B_l\).  Consequently, by
\eqref{co:slow-QH},
\begin{equation}\label{co:B-vector-bound}
 |B_l\xi|_{Q_j}^2\le C2^{2a(l-j)}B_l[\xi],\qquad j\le l.
\end{equation}
For \(j<i\), set \(m=i-j\).  The identity
\(a_i-a_j=\sum_{l=j}^{i-1}B_l\xi\) and Cauchy--Schwarz give
\[
 |a_i-a_j|_{Q_j}^2
 \le C m2^{2am}\sum_{l=j}^{i-1}B_l[\xi].
\]
Using \(N_j\le\varepsilon_\nu Q_j\), it follows that
\begin{align}
 &\sum_{i<k}\sum_{j\le i}2^{-\gamma(i-j)}N_j[a_i]
       \label{co:convolution}\\[-2mm]
 &\qquad\le C\sum_{j<k}N_j[a_j]
       +C\varepsilon_\nu
          \sum_{i<k}\sum_{j<i}(i-j)2^{-(\gamma-2a)(i-j)}
                    \sum_{l=j}^{i-1}B_l[\xi]\notag\\
 &\qquad\le C\sum_{j<k}N_j[a_j]
             +C\varepsilon_\nu H_k[\xi].\notag
\end{align}
For the last inequality, fix \(l\).  For a specified \(m=i-j\),
there are at most \(m\) intervals \([j,i)\) of length \(m\)
containing \(l\).  Hence its coefficient is bounded by
\(C\sum_{m\ge1}m^2 2^{-(\gamma-2a)m}<\infty\).
Also \(\sum_{l<k}B_l[\xi]=(H_k-H_0)[\xi]\le H_k[\xi]\).
The identity contribution is bounded separately using
\eqref{co:slow-QH}:
\begin{equation}\label{co:background-sum}
 \sum_{i<k}2^{-\gamma i}|a_i|^2
 \le C|\xi|^2\sum_{i\ge0}2^{-(\gamma-4a)i}
 \le C|\xi|^2.
\end{equation}

The lower heat decrement with \(A=1\), together with
\eqref{co:inverse-increment-order}, gives
\[
 \sum_{i<k}N_i[a_i]\le C\sum_{i<k}B_i[\xi]
 \le C H_k[\xi].
\]
Conversely, the upper heat decrement yields
\[
 H_k[\xi]\le C|\xi|^2+C\sum_{i<k}\mathcal F_i[a_i]
 \le C|\xi|^2+C\sum_{i<k}N_i[a_i]
                         +C\varepsilon_\nu H_k[\xi].
\]
Choose the rescaling late enough to absorb the last term.  This
proves \eqref{co:central-cumulative}, since \(H_0\) is comparable
with \(\Id\).  Equations \eqref{co:convolution} and
\eqref{co:background-sum} now imply
\eqref{co:weighted-total}.  Finally the lower heat decrement for
a general fixed \(A\) bounds each enlarged-mass quadratic form
on \(a_i\) by a fixed multiple of \(B_i[\xi]\).  Summing proves
the last assertion.
\end{proof}

\begin{lemma}[Changing the mean normalization]\label{co:normalization-error}
Fix \(C_0\ge1\).  For \(d(x,y)\le C_0r_i\), put
\[
 A_{i,y}=R_i(y)-Q_i(x),\qquad
 e_{i,y}=R_i(y)^{-1}\xi-H_i(x)\xi.
\]
Then
\begin{equation}\label{co:normalization-error-pointwise}
 |e_{i,y}|_{Q_i(x)}^2
 \le C\mathcal F_i(x)[H_i(x)\xi].
\end{equation}
If at each \(i<k\) at most \(J\) such centers are chosen, then
\begin{equation}\label{co:normalization-error-sum}
 \sum_{i<k}\sum_y|e_{i,y}|_{Q_i(x)}^2
 \le C_J H_k(x)[\xi].
\end{equation}
For any positive semidefinite matrices \(C_{i,y}\) satisfying
\(C_{i,y}\le\varepsilon_\nu Q_i(x)\), their error cost is
\begin{equation}\label{co:normalization-mass-error}
 \sum_{i<k}\sum_y C_{i,y}[e_{i,y}]
 \le C_J\varepsilon_\nu H_k(x)[\xi].
\end{equation}
\end{lemma}

\begin{proof}
Taking the supremum over \(v\) with \(|v|_{Q_i}=1\) in
\eqref{co:bilinear-heat-comparison} gives
\begin{equation}\label{co:bilinear-dual-form}
 A_{i,y}H_iA_{i,y}\le C\mathcal F_i.
\end{equation}
This follows directly from
\(\sup_{|v|_{Q_i}=1}|v^Tz|^2=z^TH_i z\).
Write \(a_i=H_i\xi\).  The inverse identity is
\[
 e_{i,y}=-R_i(y)^{-1}A_{i,y}a_i.
\]
Relative comparison \(R_i(y)\asymp Q_i(x)\) implies
\(R_i(y)^{-1}Q_i(x)R_i(y)^{-1}\le C H_i(x)\).
Therefore \eqref{co:bilinear-dual-form} gives
\[
 |e_{i,y}|_{Q_i}^2
 \le C a_i^T A_{i,y}H_iA_{i,y}a_i
 \le C\mathcal F_i[a_i].
\]
This proves \eqref{co:normalization-error-pointwise}.  Sum it,
use the bound \(J\), and apply
\eqref{co:weighted-total} to obtain
\eqref{co:normalization-error-sum}.  The last assertion follows
by \(C_{i,y}[e_{i,y}]\le\varepsilon_\nu|e_{i,y}|_{Q_i}^2\).
No commutation of \(R_i(y)\) with \(Q_i(x)\) has been used.
\end{proof}

\begin{lemma}[Cumulative comparison for local mass matrices]
\label{co:cumulative-comparison}
Fix constants \(A_0,C_0,L\ge1\) and a positive integer \(J\).
At each center
\(x\in B_{A_0}(p)\) and each scale \(i\ge0\), choose at most \(J\)
centers \(y\) with \(d(x,y)\le C_0r_i\).  The choice may depend
arbitrarily on \(x,i\).  Set
\begin{equation}\label{co:D-definition}
 C_{i,y}=r_i^{2-n}M(B_{Lr_i}(y)),
 \qquad D_{i,y}=R_i(y)^{-1}C_{i,y}R_i(y)^{-1}.
\end{equation}
Give the chosen centers weights \(\theta_{i,y}\in[0,1]\), and
suppose that for every \(i\) at least one has weight one and
satisfies \(B_{Lr_i}(y)\supset B_{r_i}(x)\).  Write
\begin{equation}\label{co:ZS-definition}
 Z_i=\sum_yD_{i,y},\qquad
 S_i=\sum_y\theta_{i,y}D_{i,y}.
\end{equation}
There is \(C_*<\infty\), depending only on \(n,A_0,C_0,L,J\),
such that for all sufficiently late rescalings, every such choice,
and every integer \(k\ge0\),
\begin{equation}\label{co:cumulative-estimates}
 Z_i\le o(1)H_i(x),\qquad
 \Id+\sum_{i<k}Z_i\le C_*H_k(x),\qquad
 H_k(x)\le C_*\left(\Id+\sum_{i<k}S_i\right).
\end{equation}
The \(o(1)\) is uniform over the stated centers, all scales, and
all choices satisfying these conditions.
\end{lemma}

\begin{proof}
Fix \(x\) and \(\xi\), and suppress \(x\) from the notation.
By \eqref{co:small-mass} and nearby Gram comparison,
\(C_{i,y}\le\varepsilon_\nu Q_i\), uniformly in all the choices.
Conjugation by \(R_i(y)^{-1}\), followed by relative comparison,
gives \(D_{i,y}\le C\varepsilon_\nu H_i\).  Bounded overlap
proves the first assertion of
\eqref{co:cumulative-estimates}.

For the upper sum, use
\(R_i(y)^{-1}\xi=a_i+e_{i,y}\).  Positivity gives
\[
 D_{i,y}[\xi]=C_{i,y}[a_i+e_{i,y}]
 \le2C_{i,y}[a_i]+2C_{i,y}[e_{i,y}].
\]
The ball defining \(C_{i,y}\) is contained in
\(B_{(L+C_0)r_i}(x)\).  The enlarged upper sum assertion of
Lemma~\ref{co:inverse-increments} therefore bounds the sum of
its first terms by \(C_J H_k[\xi]\).
Equation \eqref{co:normalization-mass-error} bounds the sum of
its second terms by \(C_J\varepsilon_\nu H_k[\xi]\).
Since \(\Id\le C H_0\le C H_k\), this proves the middle
assertion in \eqref{co:cumulative-estimates}.

For the reverse direction, at each scale select one of the
weight-one centers whose ball covers \(B_{r_i}(x)\).  Then
\(N_i\le C_{i,y}\).  Write
\(a_i=R_i(y)^{-1}\xi-e_{i,y}\) to obtain
\[
 N_i[a_i]\le2C_{i,y}[R_i(y)^{-1}\xi]
                +2C_{i,y}[e_{i,y}]
 \le2S_i[\xi]+2C_{i,y}[e_{i,y}].
\]
The central comparison \eqref{co:central-cumulative} and the
normalization-error bound give
\[
 H_k[\xi]\le C\left(|\xi|^2+\sum_{i<k}S_i[\xi]\right)
                      +C\varepsilon_\nu H_k[\xi].
\]
Choose the rescaling sufficiently late after all fixed constants
and absorb the last term.  The result is uniform in \(k,x,\xi\)
and the center choices.  Taking all vectors \(\xi\) proves the
last matrix inequality and completes the lemma.
\end{proof}
